% Explore, Compile, Replay
%
% Two builds from the same source (see build.sh):
%   main.pdf       named version
%   anonymous.pdf  anonymous, in review format, for a double-blind submission
%
% The anonymous build is made with \def\ECRANON{1} before this file is read.
\ifdefined\ECRANON
  \documentclass[sigconf,review,anonymous]{acmart}
\else
  \documentclass[sigconf,nonacm]{acmart}
\fi

\usepackage{tikz}
\usetikzlibrary{arrows.meta,positioning}
\usepackage{listings}
\usepackage{upquote}
\usepackage{array}
\newcolumntype{L}[1]{>{\raggedright\arraybackslash}p{#1}}

% ---- Words that change between the two builds
% The mobile tool of the case study has the same owner as ECR, so the anonymous
% build cites it without its name or its address.
\ifdefined\ECRANON
  \newcommand{\ecrartifact}{an anonymised copy of the repository, linked from the submission form}
  \newcommand{\matcite}{\cite{matanon}}
  \newcommand{\matnamed}{}
  \newcommand{\matlimit}{no two runs follow the same path, and they give the same kinds of defects, not the same list}
\else
  \newcommand{\ecrartifact}{\url{https://github.com/pierridotite/explore-compile-replay}}
  \newcommand{\matcite}{\cite{mat}}
  \newcommand{\matnamed}{, mobile-agent-testers}
  \newcommand{\matlimit}{``two runs never take the same path: expect the same kinds of bugs, not the same list''}
\fi
% The commit at which every fact about the code was checked (see claims.md).
\ifdefined\ECRANON
  \newcommand{\ecrcommit}{a pinned commit}
\else
  \newcommand{\ecrcommit}{commit \texttt{c35dfed}}
\fi

\newcommand{\code}[1]{\texttt{#1}}
\newcommand{\law}[1]{\textsf{#1}}

\lstdefinestyle{ecr}{
  basicstyle=\ttfamily\scriptsize,
  columns=fullflexible,
  keepspaces=true,
  breaklines=true,
  showstringspaces=false,
  upquote=true,
  frame=single,
  framerule=0.3pt,
  rulecolor=\color{black!35},
  xleftmargin=3pt,
  xrightmargin=3pt,
  aboveskip=3pt,
  belowskip=1pt
}
\lstset{style=ecr}

\ifdefined\ECRANON
  \setcopyright{acmlicensed}
  \acmConference[RESHAPE '27]{1st International Workshop on Reimagining Software Engineering Processes in the Agent Era}{July 2027}{Shenzhen, China}
  \acmYear{2027}
  \acmDOI{XXXXXXX.XXXXXXX}
  \acmISBN{978-1-4503-XXXX-X/2027/07}
\fi

\begin{document}

\title{Explore, Compile, Replay: Turning Persona-Driven Agent Explorations into Deterministic Tests}

\author{Pierre Raffalli}
\affiliation{%
  \city{Paris}
  \country{France}}
\email{praffalli1@gmail.com}

% Generated by scripts/tex_numbers.py from paper/results/. Do not edit.
\newcommand{\HANDtraces}{68}
\newcommand{\DSIXgoal}{12 of 15}
\newcommand{\BEprice}{8\textsuperscript{th}}
\newcommand{\BEtime}{76\textsuperscript{th}}
\newcommand{\BEpriceX}{14\textsuperscript{th}}
\newcommand{\BEtimeX}{88\textsuperscript{th}}
\newcommand{\BEreprice}{\$1.17}
\newcommand{\BEresec}{66~min}
\newcommand{\BEunrepaired}{43}
\newcommand{\BErepairprice}{\$8.22}
\newcommand{\BEdfournone}{90}
\newcommand{\BEdfourattempts}{102}
\newcommand{\CellDfour}{0 of 34 traces; no action proposed in 90 of 102 attempts}
\newcommand{\CellDsix}{0 of 102 attempts; no false green}

\begin{abstract}
Agents driven by large language models can explore an application the way a given user would, and report real defects. But an exploration is not a test: every run calls a model again, takes another path and may reach another verdict, so its findings cannot gate a merge. We present Explore-Compile-Replay (ECR), a pattern that separates three moments. An agent explores once, and the exploration is recorded as a typed trace in which actions and oracles are data. A pure compiler turns the trace into a standard test file, such as a Maestro flow, that runs in continuous integration with no model and without ECR installed. When the system changes and a replay breaks on a step, an agent or a person proposes a new action for that step, and a guard accepts it only if the change is local and every oracle gives what the trace recorded. A failed oracle is never repaired: a person decides. A defect found while exploring is recorded as an oracle that was observed to fail, and compiles to a test that stays red until the defect is fixed. An adapter contract and seven laws, checked by an executable conformance kit, state when support for a new kind of system is correct. We describe an open implementation with a reference adapter and a mobile adapter, and an evaluation protocol, written before any measured run, that injects six classes of drift, one of which is a regression that must stay red. On LocalSend, 7 traces, one of them unstable, met 25 mutants of five classes: no regression ended green, every reworded text and added dialog was detected, three of five moves and every delay went unseen, and a one-step repair was accepted for 25 of the 26 traces broken by a reworded text. Changes that need two steps were refused, and fewer than half of the text and structure mutants were repaired in full.
\end{abstract}


\begin{CCSXML}
<ccs2012>
 <concept>
  <concept_id>10011007.10011074.10011099.10011102.10011103</concept_id>
  <concept_desc>Software and its engineering~Software testing and debugging</concept_desc>
  <concept_significance>500</concept_significance>
 </concept>
 <concept>
  <concept_id>10011007.10011074.10011111.10011113</concept_id>
  <concept_desc>Software and its engineering~Software evolution</concept_desc>
  <concept_significance>300</concept_significance>
 </concept>
</ccs2012>
\end{CCSXML}
\ccsdesc[500]{Software and its engineering~Software testing and debugging}
\ccsdesc[300]{Software and its engineering~Software evolution}

\keywords{LLM agents, exploratory testing, record and replay, test oracles, test repair, continuous integration}

\maketitle


\section{Introduction}
\label{sec:introduction}

An agent driven by a large language model (LLM) can operate an application through its interface with a goal and no script~\cite{liu2023gptdroid,yoon2024droidagent,wang2024survey}. Given a persona, such as a retired teacher who reads with large text, it goes where that user would go, including screens and settings that no script was written for. Simulated users of this kind are already used to rehearse usability studies and A/B tests~\cite{lu2025uxagent,lu2025agentab}. One open-source tool runs such agents on mobile simulators in continuous integration (CI)~\matcite. In a public run of that tool against the LocalSend application, 20 agents, each on its own simulated device, filed 103 reports, grouped by hand into 46 distinct defects, in 69 minutes and for \$4.62 of model usage.

What such a run produces is a report, and a report cannot gate a merge. The documentation of the same tool states the limit: \matlimit, and it advises keeping scripted tests for regressions~\matcite. Three properties of an agent run stand in the way. It calls a model at every step, so each verdict has a delay and a price, a cost that evaluations of agents often leave out~\cite{kapoor2025agents}. It is not repeatable: the same model with the same instructions can act differently from one run to the next~\cite{ouyang2025nondeterminism,yao2025taubench}. And whether a check passes rests on the judgment of the model during the run. For a CI pipeline this is a flaky test~\cite{luo2014flaky} that also sends a bill. So the work is done twice: the agent finds a defect, then a person writes by hand the test that keeps it fixed.

Three lines of work shorten this loop. \emph{Caches} replay what the agent did last time and call the model again when the replay breaks~\cite{stagehand,midscene,shortest,arbigent,workflowuse,musclemem,skyvern}. \emph{Compilers} turn agent traces or task specifications into code or workflows, to cut the cost of inference~\cite{jaber2026auto,elyadouni2026tracecompiler,chundru2026agentic,trooskens2026compiled}. \emph{Healers} let a model repair a failing test~\cite{playwrightagents,xu2023guiding}. Each keeps something that a merge gate cannot accept (Section~\ref{sec:related}). In most of these systems, the replayed object lives inside the runtime of the tool or in the cache of a vendor, and the absence of a model during replay is a hit rate, not a guarantee. And the checks are skipped, judged again by a model, or open to rewriting by a healer: a case study of autonomous test repair reports agents that weaken assertions and delete test cases to reach a green run~\cite{lee2026limits}.

We take the position that a model belongs where judgment is needed, in choosing what to try and in proposing how to get past a button that moved, and nowhere else: not in the run that gates a merge, and not in the decision that a check may change. Explore-Compile-Replay (ECR) turns this position into a pattern with three moments (Figure~\ref{fig:loop}). An agent \emph{explores} once, as a persona, and the exploration is recorded as a typed trace in which actions and oracles are data. A pure function \emph{compiles} the trace to a file in an existing test format, which can be read, reviewed and run without ECR. CI \emph{replays} that file with no model. When the system drifts and a replay breaks on a step, a \emph{repair} is a new action for that one step, proposed by an agent or a person, and accepted by a guard only if every oracle still gives what the trace recorded. When an oracle fails, nothing is repaired: it is a regression or an intended change, and a person decides.

This paper makes four contributions.
\begin{itemize}
  \item The ECR pattern and its model: traces whose oracles are data, typed divergences, four replay outcomes that tell a recorded defect from a regression, and a bounded repair that cannot alter what a test checks (Section~\ref{sec:pattern}).
  \item An adapter contract and seven laws, with an executable conformance kit that reports for each law whether it passes, passes in part, or is skipped (Section~\ref{sec:contract}).
  \item An open implementation: a core, a command line, a reference adapter and a mobile adapter that compiles traces to Maestro flows (Section~\ref{sec:implementation}), connected to a tool that tests mobile applications with agents (Section~\ref{sec:case}).
  \item An evaluation protocol, written before any measured run, with six classes of injected drift, one of which is a regression that must stay red, and its results on LocalSend (Section~\ref{sec:evaluation}).
\end{itemize}

\section{Background and Related Work}
\label{sec:related}

\begin{table*}[t]
\caption{Systems that reuse the work of an agent or of a model instead of calling it at every step, as described by their own sources at the versions cited. ``n/s'': not stated in the cited source.}
\label{tab:related}
\footnotesize
\begin{tabular}{@{}L{2.2cm}L{3.0cm}L{3.45cm}L{3.45cm}L{3.95cm}@{}}
\toprule
System & What is replayed & A model runs during replay & Checks during replay & When the replay breaks \\
\midrule
Playwright test agents~\cite{playwrightagents} & A Playwright test file & No & Assertions in the test code & A healer agent edits the test until it passes; its instructions include fixing ``assertions and expected values'' \\
Stagehand~\cite{stagehand} & Results cached on the servers of the vendor & On a cache miss, or when the recorded selector no longer resolves & n/s & It ``falls back to full inference'' \\
Midscene~\cite{midscene} & Cached plans and element paths & On an invalid entry, and for every assertion & By a model: assertion results ``will never be cached'' & The model plans again for that run \\
Shortest~\cite{shortest} & Cached steps & On any cache error & Not evaluated again: the test is marked passed once the steps replay & A full run of the agent \\
Arbigent~\cite{arbigent} & The actions of the last passing run & For image assertions, and for a task that drifted & Image assertions, by a model & The agent runs that task again, also when an assertion fails \\
Skyvern~\cite{skyvern} & Python scripts generated from runs & When a block has no script & n/s & The agent runs instead of the script \\
Auto~\cite{jaber2026auto} & A signed WebAssembly module, with its loader & When a runtime guard fails & Differential replay against a declared contract, before the module is emitted & The reference agent takes over \\
Agentic Compilation~\cite{chundru2026agentic} & A JSON blueprint, run by its runtime & When a selector is missing or a schema is violated & n/s & A model plans again \\
Compiled AI~\cite{trooskens2026compiled} & Code generated from a specification & Not in the control plane & Validation in four stages, when the code is generated & n/s \\
\midrule
ECR (this paper) & A file in an existing test format & No, by construction: a law of the contract, checked with the network closed & Every recorded oracle, compiled to assertions of the test format & On a step: one proposed action for that step, accepted only if every oracle gives what was recorded. On an oracle: a person decides \\
\bottomrule
\end{tabular}
\end{table*}

\textbf{Oracles and replays that break.} A test needs an oracle, a way to tell correct behaviour from incorrect behaviour, and oracles remain the hard part of test automation~\cite{barr2015oracle}. Record and replay removes the writing of actions, not of oracles. Tools replay recorded streams of events~\cite{gomez2013reran}. A study of such tools in industrial cases reports where they fall short~\cite{lam2017record}, and a study of recorded web tests classifies why they break as the application changes~\cite{hammoudi2016why}. Synthesis from demonstrations generalises a recording into a program~\cite{dong2022webrobot}. ECR replaces the human demonstration by the exploration of an agent, records oracles together with actions, and gives a type to each way a replay can break.

\textbf{Exploring with agents.} LLM agents test mobile applications by conversing with the interface~\cite{liu2023gptdroid} or by setting themselves tasks~\cite{yoon2024droidagent}, and agents with personas stand in for the participants of usability tests~\cite{lu2025uxagent} and A/B tests~\cite{lu2025agentab}. Their evaluations measure what an exploration covers or finds. ECR is concerned with what remains afterwards: it takes this kind of exploration as its input and turns it into a test.

\textbf{Caching what the agent did.} Several tools record the actions of a successful run and replay them to save model calls (Table~\ref{tab:related}). In all of them the replay is an optimisation of the agent, with a path back to the model: a cache miss, an element that moved or a failed step hands control to inference again~\cite{stagehand,midscene,shortest,arbigent,workflowuse,musclemem,skyvern}. This suits automation, where the goal is to finish the task, but not a test, for two reasons. A run that should be red can come back green because the model found another way. And the checks do not survive the cache: Midscene never caches assertions, so each one calls a model~\cite{midscene}; Shortest marks a test as passed when its cached steps replay, without evaluating its expectations again~\cite{shortest}. ECR makes the replay a different program, a compiled file with no path back to a model.

\textbf{Compiling traces and specifications.} A second group compiles agent behaviour to cut inference cost in production. Auto compiles recorded spans of an agent into a signed WebAssembly module that is checked by differential replay, with runtime guards that hand control back to the agent~\cite{jaber2026auto}. TraceCompiler mines tool-call traces into a workflow representation that keeps residual LLM nodes; its authors state that it has no runtime yet and make no conformance claim~\cite{elyadouni2026tracecompiler}. Agentic Compilation compiles a page and an intent into a blueprint that its runtime executes with no model call, and plans again with a model when a selector is missing~\cite{chundru2026agentic}. Compiled AI generates code from a specification and validates it before use~\cite{trooskens2026compiled}. ECR shares the idea of compiling. It differs in what is compiled and to what: the input is an exploration with its oracles, and the output is a test in a format that teams already run and review. Outside testing, Remotion shows the same split: a video is React code that maps a frame number to an image, written once, by a person or a model, and rendered as often as needed, with seeded randomness so that parallel renders agree~\cite{remotion}. ECR takes this split as an analogy.

\textbf{Skills and experience.} Agents also keep what they learn as reusable skills or workflows~\cite{wang2024voyager,wang2025awm,zheng2025skillweaver}, as summarised experience replayed under check functions~\cite{feng2025agentrr}, or as cached plan templates~\cite{zhang2025plancaching}. In these systems the stored knowledge serves a later run of the agent, and the model stays in the loop. ECR stores a trace so that the later run needs no model.

\textbf{Repairing tests.} Repair of broken tests predates LLMs: ReAssert suggests repairs for failing unit tests~\cite{daniel2009reassert}, and WATER and VISTA repair web test scripts after the interface has changed~\cite{choudhary2011water,stocco2018vista}. LLMs now extend these techniques~\cite{xu2023guiding}. The healer agent shipped with Playwright is instructed to fix ``assertions and expected values'', among other things, and to continue ``until the test runs successfully''~\cite{playwrightagents}. A multi-agent case study reports assertion weakening and the deletion of test cases as the ways agents reached a green run, and recommends treating any change to an assertion as high risk~\cite{lee2026limits}. ECR does not ask the repairing agent to be careful. It gives it a patch format that cannot name an oracle.

What we claim is the combination, not any single element: exploration guided by a persona; oracles recorded as data; an artifact in an existing test format that runs without the tool; no model during replay, by construction and checked; a repair that is local and cannot alter an oracle; and one contract, with a conformance kit, meant for more than one kind of system.

\section{The ECR Pattern}
\label{sec:pattern}

\begin{figure*}[t]
\centering
\begin{tikzpicture}[
    font=\sffamily\footnotesize,
    >={Stealth[length=1.8mm]},
    box/.style={draw, rounded corners=1.5pt, align=center, text width=22mm, minimum height=9.5mm, inner sep=2pt},
    data/.style={box, fill=black!8},
    prog/.style={box},
    llm/.style={box, densely dashed, line width=0.8pt},
    lab/.style={font=\sffamily\scriptsize, align=center}]
  \node[data] (persona)  at (0,0)     {Persona\\[-1pt]{\scriptsize profile, goal}};
  \node[llm]  (explore)  at (2.95,0)  {Explore\\[-1pt]{\scriptsize an agent, once}};
  \node[data] (trace)    at (5.9,0)   {Trace\\[-1pt]{\scriptsize actions and oracles}};
  \node[prog] (compile)  at (8.85,0)  {Compile\\[-1pt]{\scriptsize pure function}};
  \node[data] (artifact) at (11.8,0)  {Artifact\\[-1pt]{\scriptsize standard test file}};
  \node[prog] (replay)   at (14.75,0) {Replay in CI\\[-1pt]{\scriptsize no model}};
  \draw[->] (persona) -- (explore);
  \draw[->] (explore) -- (trace);
  \draw[->] (trace) -- (compile);
  \draw[->] (compile) -- (artifact);
  \draw[->] (artifact) -- (replay);

  \node[prog] (guard)    at (5.9,-2.15)   {Repair guard\\[-1pt]{\scriptsize accepts or refuses}};
  \node[llm]  (proposal) at (11.8,-2.15)  {Proposal\\[-1pt]{\scriptsize one step, one action}};
  \node[prog] (person)   at (14.75,-2.15) {A person\\[-1pt]decides};
  \draw[->] (replay.south) -- node[lab, right] {an oracle\\failed} (person.north);
  \draw[->] ([xshift=-8mm]replay.south) -- ++(0,-0.6) -| node[lab, above, pos=0.25] {a step diverged} (proposal.north);
  \draw[->] (proposal) -- node[lab, above] {from an agent or a person} (guard);
  \draw[->] (guard) -- node[lab, right] {accepted:\\new action} (trace);

  \node[llm, text width=5mm, minimum height=3.2mm] (legend) at (-0.75,-2.15) {};
  \node[lab, right=1mm of legend, align=left] {a model may run\\here, and only here};
\end{tikzpicture}
\caption{The ECR loop. A model may run in the two dashed boxes only. Compiling and replaying are deterministic programs. When a step diverges, a proposal of one new action goes through a guard. When an oracle fails, nothing is repaired.}
\Description{A pipeline from left to right: persona, explore, trace, compile, artifact, replay in continuous integration. From replay, a step divergence leads to a proposal of one action, then to a repair guard that updates the trace if it accepts. An oracle failure leads to a person.}
\label{fig:loop}
\end{figure*}

\subsection{Traces}

A \emph{system} is anything a program can act on and observe: a mobile application on an emulator, a web page, an API. An \emph{adapter} for a kind of system fixes two typed vocabularies: \emph{actions}, such as a tap on an element, and \emph{checks}, such as ``this element is visible''. A \emph{persona} is a profile, a goal, and an optional context such as device, locale and text size.

A \emph{trace} records one exploration of one system by one persona. It holds the persona, a list of parameters, a list of invariants, an ordered list of steps, and its provenance. A \emph{step} has a stable identifier, an optional intent in plain words, one action, and its \emph{checkpoints}: oracles that must hold right after the step. An \emph{invariant} is an oracle that must hold after every step. An \emph{oracle} is an identifier, a check from the vocabulary, and what was observed during the exploration, \emph{held} or \emph{failed} (Section~\ref{sec:bug}).

Three rules make a trace safe to commit and to review. A parameter is stored by name and its value is supplied at replay, so no password reaches the repository. Evidence attached to a step, such as a screenshot or the point where the finger landed, is for people and for tools that display a trace; no compiler reads it. Intent and provenance are kept for audit and are never used at replay.

\subsection{Explore}

Exploring has two halves. Choosing the steps is the work of the agent, of a person, or of an automated testing tool that hands over to a model where it stops making progress, as LLMDroid does when coverage growth slows~\cite{wang2025llmdroid}. Recording them is the work of ECR: each action is executed through a driver, the system is observed until it reports a stable state, and each oracle is verified by a deterministic procedure of the driver that calls no model. An oracle enters the trace as a guard only if it held. If it did not hold, the recording stops with an error. It also stops when the system never settles.

\subsection{Compile and replay}

A \emph{target} is a pair of functions for one artifact format. \emph{Compile} is pure: the same trace gives the same bytes. Its result, the \emph{artifact}, is a file in an existing format, together with the SHA-256 hash of the canonical form of the trace and a \emph{source map} that gives each step one range of lines. Figure~\ref{fig:step} shows a step and the lines it compiles to. The hash lets CI refuse a compiled file that no longer matches its trace.

\begin{sloppypar}
\emph{Replay} runs the artifact with its own tool and returns a verdict with a list of typed \emph{divergences}. Four types say that a step could not be carried out: \code{target\_not\_found}, \code{action\_failed}, \code{unexpected\_state} and \code{timeout}. The fifth, \code{oracle\_failed}, says that the steps ran and that the system no longer does what it did. Everything that follows rests on this distinction: the first four may be repaired, the fifth never is.
\end{sloppypar}

\begin{figure}[t]
\begin{lstlisting}
{ "id": "s8",
  "action": { "kind": "tapOn", "target": { "text": "Save" } },
  "checkpoints": [
    { "id": "item-listed", "check": { "kind": "visible",
        "target": { "text": { "param": "ITEM_TITLE" } } } },
    { "id": "empty-message-gone", "check": { "kind": "notVisible",
        "target": { "text": "Nothing here yet (0 items)" } } } ] }
\end{lstlisting}
\begin{lstlisting}
# Step s8
- tapOn:
    text: "Save"
# Oracle item-listed
- assertVisible:
    text: "${ITEM_TITLE}"
# Oracle empty-message-gone
- assertNotVisible:
    text: "Nothing here yet \\(0 items\\)"
# Oracle no-crash-dialog
- assertNotVisible:
    text: "App keeps stopping"
\end{lstlisting}
\caption{A step of the example mobile trace of the repository (top, white space condensed, evidence left out) and the lines of the Maestro flow it compiles to (bottom). The parameter stays a name, the literal text is escaped, and the invariant of the trace is checked after the step.}
\Description{A JSON object for one step with an action and two checkpoints, and the twelve lines of YAML it compiles to.}
\label{fig:step}
\end{figure}

\subsection{A defect is a test}
\label{sec:bug}

An agent that explores finds defects. In a report, a defect is a paragraph. In a trace, it is an oracle: it states what should be true, and it is recorded as \emph{failed} because it was false during the exploration. ECR records it only if the driver confirms that it is false. It compiles like any other oracle, so the compiled test is red until the defect is fixed: a reproduction that anyone can run, with no model.

Because a trace now says what it expects, a replay has four outcomes instead of two. \emph{Pass}: every oracle holds and the trace records no defect. \emph{Known bug}: the replay stops on the first oracle recorded as failed, and on nothing before it. \emph{Fixed}: that oracle now holds, and it can be kept as a guard. \emph{Fail}: anything else, and the only outcome that fails the build. Two limits follow. A replay stops at the first recorded defect, so the steps after it are not checked until it is fixed. And only a checkpoint can record a defect, never an invariant.

\subsection{Bounded repair}
\label{sec:repair}

When a replay ends on a step divergence, someone has to look at the changed system and say how the step should now be carried out. ECR does not produce that answer. An agent or a person proposes it, and ECR decides. A \emph{proposal} is one step identifier and one action, and nothing else. The guard runs five checks in order and refuses, with a reason, at the first one that fails:
\begin{enumerate}
  \item the replay of the current trace diverges, and not on an oracle;
  \item the proposal names the step that diverged;
  \item the action belongs to the vocabulary and compiles;
  \item in the compiled file, the text of every other step is unchanged;
  \item the replay of the patched artifact gives what the trace recorded.
\end{enumerate}
Two properties follow from this design. A proposal cannot remove or weaken an oracle, because it has no way to name one. And a patch that gets past the broken step by doing something else, such as ticking another item of a list, is refused by the fifth check, because an oracle then fails. An accepted repair is one changed action in the trace file, which reads as one hunk in a review. These bounds have a price: a change that needs a new step is refused, and the path has to be explored again.

\section{The Adapter Contract and Seven Laws}
\label{sec:contract}

To support a new kind of system, one writes an \emph{adapter}: the two vocabularies as schemas, a driver, and one target for each artifact format (Figure~\ref{fig:contract}). The \emph{driver} acts on the live system. It is used to explore, and a replay never uses it: a replay runs the compiled artifact. Its \code{verify} must be deterministic and call no model. Its \code{snapshot} and \code{restore} are optional and declared as a capability. A \emph{target} compiles and replays.

\begin{figure}[t]
\begin{lstlisting}
interface SystemDriver<State, A, O> {
  capabilities: { snapshot: boolean };
  setup(): Promise<void>;  teardown(): Promise<void>;
  observe(): Promise<{ state: State; stable: boolean }>;
  execute(action: A): Promise<void>;
  verify(oracle: O): Promise<Verdict>;  // deterministic, no model
  snapshot?(): Promise<SnapshotId>;
  restore?(id: SnapshotId): Promise<void>;
}
interface Target<A, O> {
  id: string;  // "maestro", "playwright"
  compile(trace: TypedTrace<A, O>): CompiledArtifact;  // pure
  replay(artifact: CompiledArtifact, options?: ReplayOptions):
    Promise<ReplayResult>;
  replayInstrumented?(artifact: CompiledArtifact,  // for law L2
    options?: ReplayOptions):
    Promise<{ result: ReplayResult; received: Received[] }>;
}
interface Adapter<State, A, O> {
  id: string;  actions: z.ZodType<A>;  oracles: z.ZodType<O>;
  driver: SystemDriver<State, A, O>;  targets: Target<A, O>[];
}
\end{lstlisting}
\caption{The adapter contract of \code{spec/contract.ts}, abridged: comments and \code{readonly} modifiers are left out.}
\Description{Three TypeScript interfaces: SystemDriver, Target and Adapter.}
\label{fig:contract}
\end{figure}

An adapter is correct when it obeys seven laws (Table~\ref{tab:laws}). The laws are executable. A conformance kit takes an adapter and a few traces, and reports one status for each law. The kit says what it cannot check: a law that is covered in part is reported as \emph{partial}, a law with no applicable check as \emph{skipped}, and neither counts as a pass.

\begin{table}[t]
\caption{The seven laws, and the report of the conformance kit for the reference adapter at \ecrcommit{}, in the CI job that closes the network.}
\label{tab:laws}
\footnotesize
\begin{tabular}{@{}L{1.75cm}L{5.25cm}l@{}}
\toprule
Law & Statement & Kit \\
\midrule
\law{L1} Determinism & The same trace compiles to the same bytes. The same artifact, on the same system state, gives the same verdict. & pass \\
\law{L2} Fidelity & The artifact runs the actions of the trace and checks its oracles, in the same order. It adds none and drops none. & pass \\
\law{L3} Traceability & Each step maps to one distinct range of the artifact. Ranges do not overlap. & pass \\
\law{L4} No model & Compiling and replaying call no model. & pass \\
\law{L5} State & If the driver declares snapshots, restoring a snapshot gives a state that is observed as identical. & pass \\
\law{L6} Safe repair & A patch stays inside the range of the step that diverged. It removes no oracle and weakens none. A failed oracle is never repaired. & partial \\
\law{L7} No secret & Neither the trace nor the artifact contains a secret value. & pass \\
\bottomrule
\end{tabular}
\end{table}

Table~\ref{tab:laws} gives the report for the reference adapter of Section~\ref{sec:implementation}. Six laws pass. For \law{L1} the kit compiles each trace twice and compares the bytes, then replays each artifact three times and compares the results. For \law{L2} a recorded trace must replay as it was observed. The artifact is then replayed against an instrumented copy of the system, which notes every call it receives, and the kit compares these notes with the trace: the same actions with the same values, one read of the state for each oracle, in the same order, and nothing else. For \law{L3} it validates the source map. For \law{L5} it observes, takes a snapshot, acts, restores, and compares the two observations. For \law{L7} it searches the trace and the artifact for the value of each parameter that the trace declares secret.

\law{L4} needs help from outside. The kit verifies that no artifact imports a known model SDK, but it does not close the network itself: the caller does, and says so. The kit then checks that no public address answers before it runs anything. A CI job runs it this way, in a network namespace that has only the loopback. Otherwise the kit reports \law{L4} as partial.

One law is partial, and the report says which part is missing. For \law{L6}, the kit gives each step the action of another step and verifies, on the compiled files, that the text of every other step is unchanged; the refusal of every patch after a behaviour change is covered by the tests of the core, not yet by the kit. The kit has been run on the reference adapter only. On the mobile adapter it needs a device, and the mobile target has no instrumented replay, so \law{L2} would stay partial there.

\section{Implementation}
\label{sec:implementation}

ECR is open source under the Apache~2.0 licence. It is written in TypeScript, runs on Node.js~22 with no build step, and has one runtime dependency, a schema library. This section describes it at \ecrcommit{}. Sizes are lines that are neither blank nor comments, by the counting rule of the repository.

\textbf{Core and command line.} The schema of a trace and the contract are two files of 85 and 57 lines. The core (337 lines) records a trace through a driver, validates it against the schema and the vocabularies of its adapter, hashes it, resolves parameters, checks a source map, classifies a replay into the four outcomes of Section~\ref{sec:bug}, and guards a repair. It imports no model SDK. The conformance kit is 260 lines. A command line exposes five commands: \code{validate}, \code{compile}, \code{replay}, \code{repair} and \code{conformance}.

\textbf{Reference adapter.} The reference adapter (302 lines) drives a to-do list held in memory and compiles a trace to a test file for the test runner built into Node.js. It implements snapshots and an instrumented copy of its system, so that every law has a subject. Two altered versions of the system serve as fixtures. In the first, an item is no longer found by its title. In the second, a new item starts ticked: for a trace recorded on the original system this is a change of behaviour, and for a trace explored on it this is a system with a defect.

\textbf{Mobile adapter.} The mobile adapter (298 lines) defines nine actions (launch, tap, long press, type, erase, key press, swipe, scroll until visible, hide the keyboard) and two checks (visible, not visible). A step finds its element by accessibility identifier, by visible text or, as a last resort, by a point given in percent of the screen. A press can say where in its element it lands, and a swipe can keep its two ends. Traces compile to Maestro flows~\cite{maestro}: a parameter becomes a variable passed on the command line, and a literal text is escaped because Maestro reads it as a regular expression. In CI, the syntax checker of Maestro validates the compiled flows. Replay runs Maestro. When a flow fails, the target finds the command that failed, in the record that Maestro keeps or in what it prints, and maps it back through the source map: a failed assertion becomes \code{oracle\_failed} with the oracle it checks, and a failed action becomes a divergence on its step, whose type is read from the wording of the message of Maestro. That reading is the weakest part of the target: on a device it has been checked for one cause only, an element that is not found (Sections~\ref{sec:evaluation} and~\ref{sec:case}), and it never reports \code{unexpected\_state}. The adapter has no driver: an external tool explores and writes the trace (Section~\ref{sec:case}).

\textbf{On a device.} A CI job, started by hand, builds the demonstration application of Section~\ref{sec:case}, starts an Android~11 emulator and replays the flows compiled from two traces, with no model. Both traces were written directly from the source, not explored by an agent. In its run of 25~November 2026, the first replay passed. The second stopped on the oracle that its trace records as failed, a counter that reads ``1~tasks'', and ECR reported it as a known bug. It shows that the chain works on a device.

\textbf{Tests, and what is not written.} A suite of 69 tests covers the packages. Two parts are not started: a server through which an agent explores while ECR records each step, and a web adapter with a Playwright target.

\section{Evaluation}
\label{sec:evaluation}

The protocol of this evaluation was written before any measured run and is part of the artifact, with six dated amendments. As guidelines for empirical studies that involve LLMs ask~\cite{baltes2025guidelines}, the agent's model version, prompts and traces are published, and every number below comes from a result file.

\subsection{Questions}

\textbf{RQ1, cost.} What does a verdict cost with ECR, against running the agent again, and after how many runs does one exploration plus its replays and repairs cost less? \textbf{RQ2, determinism.} On an unchanged system, does a replay give the same result 20 times of 20, and does the agent, run again five times? \textbf{RQ3, drift.} When the system changes, does the replay break where it should, is a correct repair accepted and a wrong one refused, and does a regression stay red? \textbf{RQ4, generality.} What does a new adapter cost to write, and do the seven laws transfer to it?

\subsection{Injected drift}

Real drift arrives slowly and without labels, so we inject it, in the manner of mutation analysis~\cite{jia2011mutation}. A \emph{mutant} is one change to the system under test, applied alone, written once the traces are frozen and committed before any replay on a mutant. Table~\ref{tab:drift} gives the six classes, what the protocol expects of each, and what we observed. A change of behaviour (\textsf{D6}) is the control: a \textsf{D6} mutant that ends green, by any path, counts as a false green and as a failure of law \law{L6}.

\begin{table*}[t]
\caption{Injected drift on LocalSend, five mutants per class: what the protocol expects (columns 3 and 5), and what 3 replays of the 7 traces per mutant and at most 3 proposals per diverged trace gave. Detected: a trace fails its 3 replays.}
\label{tab:drift}
\footnotesize
\begin{tabular}{@{}L{1.55cm}L{2.85cm}L{2.9cm}L{2.05cm}L{1.15cm}L{5.05cm}@{}}
\toprule
Class & The change & Expected replay & Detected (type) & Repair & Repairs accepted \\
\midrule
\textsf{D1} Text & A targeted text is reworded & \code{target\_not\_found} & 5 of 5 (5 of 5) & Accepted & 25 of the 26 traces; 4 of 5 mutants in full \\
\textsf{D2} Identifier & A targeted identifier is renamed & \code{target\_not\_found} & \multicolumn{3}{@{}l}{Not applicable: the 54 taps of the 7 traces all target a text} \\
\textsf{D3} Structure & The target moves & \code{target\_not\_found}, \code{action\_failed} & 2 of 5 (1 of 2) & Accepted & 0 of 8 traces: 2 steps to change (7), oracle (1) \\
\textsf{D4} Added step & A dialog comes first & \code{target\_not\_found}, \code{unexpected\_state} & 5 of 5 (0 of 5) & Refused & \CellDfour \\
\textsf{D5} Delay & The target comes 1, 5 or 15~s later & Pass, or \code{timeout} & 0 of 5; 102 of 105 replays pass & None & None asked \\
\textsf{D6} Behaviour & A checked state is not reached & \code{oracle\_failed} & 5 of 5 (5 of 5) & Refused & \CellDsix \\
\bottomrule
\end{tabular}
\end{table*}

\subsection{Procedure and baselines}

The subject is LocalSend~v1.19.0, an open-source file-sharing app that the author did not write, built from source and run on Android emulators; the agent runs the model \code{claude-sonnet-5-5}. A trace keeps the path of one exploration and the checks that an AI assistant wrote on it, under a rule fixed in advance. Each trace is replayed on the unchanged app, then three times on each mutant. For each trace that diverges, the exploring agent, with a fixed prompt, proposes a new action at most three times; it sees the steps without their oracles, the divergence and the current screen, and may answer that one action is not enough. The guard decides; an AI assistant judges accepted repairs.

Two baselines run on the same traces: the same agent run again (RQ1, RQ2, and three runs on each \textsf{D6} mutant), and a flow that an AI assistant keeps green by editing the trace for each mutant on which it stops (lines changed in the compiled flow). The protocol lists in advance the outcomes that count against ECR.

\subsection{Results}

\textbf{Explorations and traces.} Twenty explorations, two per persona, cost a median of \$0.20 (0.16 to 0.36), 23.5 turns and 112.5~s of agent time. Replayed once, 10 of the 20 paths passed (Wilson 95\%: 0.30 to 0.70); with their checks, 7 traces from 4 personas, fewer than the 10 planned, passed and were kept (62 checks, none recording a defect).

\textbf{RQ1, cost.} A replay calls no model (the last replay of each trace ran with the network closed, and passed) and takes 52 to 101~s per trace. An agent run of the same persona costs a median of \$0.20 (0.12 to 0.31) and 121.5~s. A replay thus saves the whole price of a run but only about a minute. With every repair asked in the campaign, \textsf{D6} included (\BEreprice{} and \BEresec{} of device time per trace), a kept trace costs less than the agent from the \BEprice{} run in price and the \BEtime{} in time; from the \BEpriceX{} and the \BEtimeX{} if each of the \BEunrepaired{} pairs of a trace and a mutant left unrepaired outside \textsf{D6} costs a new exploration. How long a trace lives before it drifts is not measured.

\textbf{RQ2, determinism.} On the unchanged app, 137 of 140 replays gave what the trace records: six traces 20 times of 20, one 17 times, failing each time on the check right after a Back (cause not established), and once more in the 21 replays of the unpatched build that preceded the mutants. This counts against ECR. Run again, the agent reached the goal of a trace, the screen that its last checks name, in 30 of 35 trace-runs (Wilson: 0.71 to 0.94), never for the unstable trace, and its own path replayed in full in 12 of 20 runs.

\textbf{RQ3, drift.} No \textsf{D6} mutant ended green (Table~\ref{tab:drift}). Every reworded text stopped the replay with \code{target\_not\_found}. Every added dialog was detected, but mostly as a failed oracle, which is never repaired: a wrong type, on the safe side. Three moves and every delay went unseen: a text target survives a move, and the waits of Maestro absorb up to 15~s; still, a 5~s delay failed one trace on an oracle in 2 replays of 3. The 25 accepted proposals are all on \textsf{D1}: 22 taps by text and 3 by position, 21 at the first attempt; an attempt took a median of one turn and \$0.03. An AI assistant judged all 25 right. Each kept the hashes of the oracles, and 24 then passed 3 replays of 3, the other being on the unstable trace. The tab swap of \textsf{D3} needed two steps changed in every trace, as did the one \textsf{D1} trace that taps the reworded button twice, and all their proposals were refused: the bound of one step. Of the 7 mutants of \textsf{D1} and \textsf{D3} that stop a trace, 4 were repaired for every trace, and 4 of the 10 in all: fewer than half, which counts against ECR. On \textsf{D4} the proposer itself declined in \BEdfournone{} of \BEdfourattempts{} attempts. Edited freely, the same \HANDtraces{} broken traces needed 2 to 4 lines changed in the compiled flow, \textsf{D3} and \textsf{D4} included, and all then replayed 3 times of 3 apart from the unstable trace: free edits reach what the bound refuses, with no guard on what they change. Run again on each \textsf{D6} mutant, the agent reached its goal screen in \DSIXgoal{} runs; its findings were not read.

\textbf{RQ4, generality.} The reference adapter has 302 lines and the mobile adapter 298, with no driver and no change to the contract. The kit runs on the reference adapter; read from its tests and this campaign, the mobile adapter passes \law{L3} and \law{L4}, holds \law{L1} and part of \law{L2} for 6 traces of 7 and \law{L6} on every accepted repair and \textsf{D6} mutant; \law{L5} and \law{L7} are not checked.

\section{Case Study: Agents That Test Mobile Applications}
\label{sec:case}

The mobile adapter was built for an existing open-source tool\matnamed~\matcite. The tool starts one coding agent per device on iOS Simulators and Android emulators in CI. Each agent plays a persona, drives the application with no script, and writes a report.

Before ECR, the path that led an agent to a defect survived only as prose in its report, and the documentation advised writing a Maestro flow by hand for each defect worth guarding. Three pieces now connect the tool to ECR, and the trace format is the only thing they share: a recorder, a map, and the Maestro target of Section~\ref{sec:implementation}.

\textbf{A recorder.} Each Android agent leaves a trace next to its report. Nothing changes for the agent: a pass-through in front of the device bridge notes its taps, and a script writes the trace when the agent is done. A tap is stored by the name of the element that was under the finger, failing that by its identifier, and by a position only as a last resort. Every text that the agent typed becomes a parameter, and the values go to a separate file, not to be committed.

\textbf{A map.} A trace written by an agent is a path, not yet a test: it has no oracle. A local web page draws the screens of the application as a map, with the paths of the agents, a heat map of their taps and their findings. On a screen, a person picks an element and says whether it must be visible or must not be: the check becomes an oracle of the step that led to that screen. Started from a finding, the same gesture records the check as failed, as in Section~\ref{sec:bug}. Only checks can change: the page refuses a trace whose steps differ from the explored ones. No driver verifies these checks when they are written. The first replay on the unchanged application does, and a trace is kept only if that replay gives what the trace records.

\textbf{State.} The recorder and the map are on the main branch of the tool, for Android only, and the map has not been tried on a real run. The tool also replays the path of an agent right after its exploration, with no model key and no edit in between, and holds the proposer of repairs used in Section~\ref{sec:evaluation}. Trials on its demonstration application, then on LocalSend, found errors of the recorder, such as a button and its label counted as two elements where Maestro finds one, that were fixed before the measured runs.

\section{Limitations and Threats to Validity}
\label{sec:limits}

\textbf{What a deterministic oracle cannot say.} A persona notices that a label is confusing or that a screen feels slow. A check such as ``this text is visible'' cannot express that. ECR keeps such findings in the report and turns into tests only the ones that can be stated as a check. The mobile vocabulary has two checks today, which narrows this further.

\textbf{Who writes the oracles.} In the mobile case study the agent supplies a path and a finding, and a person writes the checks. In the evaluation an AI assistant writes them under a fixed rule, and also keeps the hand baseline green: RQ1 counts no cost for writing the checks, which may be weaker, and the baseline is weaker than one kept by a person who knows the app. Recording oracles while the agent explores needs the server that is not written yet.

\textbf{The price of bounds.} A repair is one action for one step. On LocalSend this bound accepted every reworded text but one, and refused every added dialog and a swap of two tabs, which needed two steps changed. Fewer than half of the mutants of \textsf{D1} to \textsf{D3} were repaired in full. A wider repair, several steps under the same guard, is the obvious next step, and it widens what a wrong repair can reach.

\textbf{A guard as strong as its oracles.} The repair guard compares outcomes, not intentions. A new action that reaches the same oracles by another route is accepted, and for five of the 25 accepted repairs no oracle could detect a missed tap. This is why each accepted repair is judged afterwards, here by an AI assistant.

\textbf{Determinism is relative to the system.} Law \law{L1} holds for the same system state. On a device, state leaks between runs and timing varies. The contract offers snapshots and a flag for unsettled states, but a flaky system gives a flaky replay. A tap stored as a point breaks on another screen size.

\textbf{Secrets at replay.} Values never reach the trace or the artifact, but the mobile target passes them to Maestro on the command line, where the process list shows them while a flow runs.

\textbf{Threats to the evaluation.} The drift is injected, not observed, and the mutants were written inside the project, by an AI assistant. Delays and moves that a text target survives went unseen, so the detection rate of \textsf{D3} and \textsf{D5} says as much about Maestro as about ECR. Seven traces from four personas are few, and ECR does not measure what a set of traces covers. We limit this with a protocol committed before any measured run, mutants committed before any replay on them, a control class that must stay red, and published result files. One app, one platform and one model say little about others. One person designed ECR and its evaluation.

\section{Implications for the Development Process}
\label{sec:implications}

\textbf{A place for the model.} ECR assigns the model to two activities, exploring and proposing, and both produce files that are reviewed before they count. Exploration becomes scheduled work whose output arrives as a trace and a compiled test in a change request, while the gate stays a deterministic program. The two decisions that change what a test means stay with people: a repair reaches a reviewer as one changed action, already checked against every recorded oracle, and a failed oracle reaches a person as a question, regression or intended change, that no agent may settle by editing the test.

\textbf{Findings as executable artifacts.} A defect found by an agent reaches its bug report as a test that anyone can run, red until the fix and a guard after it.

\textbf{Economics.} The cost of a model moves from each run to each change: one exploration when a path is first tested, one proposal when the system drifts under it. On LocalSend a replay saved the whole price of an agent run but only about a minute of its time, and repairs cost mostly device time, so ECR paid back in price within fifteen runs and in time only after tens. How often a system drifts, a property of each project, decides the rest.

\section{Conclusion}
\label{sec:conclusion}

Agents explore well and repeat badly. ECR keeps the exploration and removes the repetition: the agent explores once, a pure function compiles what it did and what it checked into a standard test, and CI replays that test with no model. A repair is a new action for one step, accepted only if every recorded oracle still gives the same result, and a failed oracle is left to a person. A defect becomes a test that is red until it is fixed. A contract and seven executable laws state what a new adapter must guarantee. On LocalSend, replays held for six traces of seven, none of the five injected regressions ended green, and the bound of one step was both safe and narrow.

\section*{Artifact Availability}

The specification, the implementation, the conformance kit, the mutants, the evaluation protocol and the result files are under the Apache~2.0 licence at \ecrartifact; the code is described at \ecrcommit{}. The measured traces and the prompts are with the mobile tool~\matcite.

\section*{Use of Generative AI}

ECR exists to bound what a model does in a test pipeline, and this work was done with the same tools it studies. An LLM coding agent wrote most of the implementation and the first draft of this text, under the direction of the author, who reviewed the code and the claims and is responsible for both. An AI assistant wrote the checks, the mutants and the hand-kept flows of the evaluation, and judged the accepted repairs, under the rules of the protocol. The model of the exploring agent is named, with its version, in the protocol and in the result files of its runs.


\bibliographystyle{ACM-Reference-Format}
\bibliography{refs}

\end{document}
